%ARQUIVO DE PREAMBULO DA TESE - PACOTES E CONFIGURA\c{C}\~{O}ES

\documentclass[
	% -- op\c{c}\~{o}es da classe memoir --
	12pt,				% tamanho da fonte
	openright,			% cap\'{\i}tulos come\c{c}am em p\'{a}g \'{\i}mpar (insere p\'{a}gina vazia caso preciso)
	oneside,			% para impress\~{a}o em verso e anverso. Oposto a oneside
	a4paper,		% tamanho do papel.
	% -- op\c{c}\~{o}es da classe abntex2 --
	%chapter=TITLE,		% t\'{\i}tulos de cap\'{\i}tulos convertidos em letras mai\'{u}sculas
	%section=TITLE,		% t\'{\i}tulos de se\c{c}\~{o}es convertidos em letras mai\'{u}sculas
	%subsection=TITLE,	% t\'{\i}tulos de subse\c{c}\~{o}es convertidos em letras mai\'{u}sculas
	%subsubsection=TITLE,% t\'{\i}tulos de subsubse\c{c}\~{o}es convertidos em letras mai\'{u}sculas
	% -- op\c{c}\~{o}es do pacote babel --
	english,			% idioma adicional para hifeniza\c{c}\~{a}o
	%french,			% idioma adicional para hifeniza\c{c}\~{a}o
	%spanish,			% idioma adicional para hifeniza\c{c}\~{a}o
	brazil,				% o \'{u}ltimo idioma \'{e} o principal do documento
	sumario=tradicional,
	]{abntex2}

% Stop \paragraph run-in headings from being numbered (X.Y.Z.0.N) or listed in
% the ToC. The body uses no \subsubsection and many \paragraph emphatic leads;
% numbering/listing down to subsection keeps the Sumario clean (X.Y.Z).
\setsecnumdepth{subsection}
\settocdepth{subsection}

% ---
% PACOTES
% ---

% ---
% Pacotes fundamentais
% ---
\usepackage{cmap}				        % Mapear caracteres especiais no PDF
\usepackage{lmodern}			      % Usa a fonte Latin Modern			
\usepackage[T1]{fontenc}		    % Selecao de codigos de fonte.
\usepackage[utf8]{inputenc}		  % Codificacao do documento (convers\~{a}o autom\'{a}tica dos acentos)
\usepackage{lastpage}			      % Usado pela Ficha catalogr\'{a}fica
\usepackage{indentfirst}		    % Indenta o primeiro par\'{a}grafo de cada se\c{c}\~{a}o.
\usepackage{color}				      % Controle das cores
\usepackage[pdftex]{graphicx}	  % Inclus\~{a}o de gr\'{a}ficos
\usepackage{epstopdf}           % Pacote que converte as figuras em eps para pdf
\usepackage{lipsum}             % Pacote que gera texto dummy
\usepackage{blindtext}          % Pacote que gera texto dummy
% ---
		
% ---
% Pacotes adicionais, usados apenas no \^{a}mbito do Modelo Can\^{o}nico do abnteX2
% ---
%\usepackage{nomencl}

\let\printglossary\relax
\let\theglossary\relax
\let\endtheglossary\relax

\usepackage[nonumberlist]{glossaries} % nonnumberlist nao mostra as paginas nas quais os acronimos aparecem no texto
\usepackage{amsmath}
\usepackage{amssymb}
\usepackage{bbm}
\usepackage[chapter]{algorithm}
\usepackage{algorithmic}
\usepackage{multirow}
\usepackage{rotating}
\usepackage{pdfpages}
\usepackage{subcaption}
\usepackage{booktabs}
\usepackage{tabularx}
\usepackage{pdflscape}
\usepackage{chngcntr}
\counterwithin{figure}{chapter}
\counterwithin{table}{chapter}

% ---
% Pacotes de cita\c{c}\~{o}es
% ---
\usepackage[brazilian,hyperpageref]{backref}	 % Paginas com as cita\c{c}\~{o}es na bibl
\usepackage[alf,abnt-etal-cite=2,abnt-etal-list=0,abnt-etal-text=emph]{abntex2cite}	% Cita\c{c}\~{o}es padr\~{a}o ABNT
% ---
% Pacote de customiza\c{c}\~{a}o - Unicamp
% ---
\usepackage{unicamp}
% ---
% CONFIGURA\c{C}\~{O}ES DE PACOTES
% ---

% ---
% Configura\c{c}\~{o}es do pacote backref
% Usado sem a op\c{c}\~{a}o hyperpageref de backref
\graphicspath{{eps/}{figs/}{cap1/}}
\DeclareGraphicsExtensions{.pdf,.png,.eps,.jpg}

% Line-breaking relief: lets TeX absorb the few minor overfull boxes caused by
% unbreakable long \texttt{} tokens (file paths, class names) without loosening
% body text or altering any wording.
\emergencystretch=3em
\tolerance=1200

%customiza\c{c}\~{a}o do negrito em ambientes matem\'{a}ticos
\newcommand{\mb}[1]{\mathbf{#1}}
\newcommand{\abs}[1]{\left|#1\right|}
\newcommand{\norm}[1]{\left\|#1\right\|}
\newcommand{\partialorder}{\cdot \geq}
\newcommand{\h}{\mathrm{h}}
\newcommand{\x}{\mathrm{x}}
\newcommand{\z}{\mathrm{z}}
\newcommand{\entropia}[1]{H{(#1)}}
%customiza\c{c}\~{a}o de teoremas
\newtheorem{mydef}{Defini\c{c}\~{a}o}[chapter]
\newtheorem{lemm}{Lema}[chapter]
\newtheorem{theorem}{Teorema}[chapter]
\floatname{algorithm}{Pseudoc\'{o}digo}
\renewcommand{\listalgorithmname}{Lista de Pseudoc\'{o}digos}
\newenvironment{proof}[1][Prova]{\begin{trivlist}
\item[\hskip \labelsep {\bfseries #1}]}{\end{trivlist}}
\newcommand{\qed}{\hfill \ensuremath{\Box}}

% --- Auto-generated thesis macros (numbers & tables) ---
% Generated by scripts/thesis/render_tables.py — do not edit by hand.
% Auto-generated from docs/results/ablation/Falpha_summary.json and docs/results/ablation/Fcoupling_summary.json -- do not hand-edit.
% Regenerate with: PYTHONPATH=. .venv/bin/python -m scripts.thesis.render_tables
\newcommand{\NumSeeds}{10}
\newcommand{\AlphaZeroPPO}{+138.6}
\newcommand{\AlphaZeroPPOCILow}{+129.0}
\newcommand{\AlphaZeroPPOCIHigh}{+147.7}
\newcommand{\AlphaZeroDQN}{+116.6}
\newcommand{\AlphaZeroDQNCILow}{+102.6}
\newcommand{\AlphaZeroDQNCIHigh}{+130.7}
\newcommand{\AlphaZeroATwoC}{+147.1}
\newcommand{\AlphaZeroATwoCCILow}{+137.9}
\newcommand{\AlphaZeroATwoCCIHigh}{+156.2}
\newcommand{\AlphaZeroRF}{+136.5}
\newcommand{\AlphaZeroOracle}{+194.8}
\newcommand{\AlphaZeroGap}{+2.1}
\newcommand{\AlphaZeroVerdict}{overlapping}
\newcommand{\AlphaTwoPPO}{+121.6}
\newcommand{\AlphaTwoPPOCILow}{+109.1}
\newcommand{\AlphaTwoPPOCIHigh}{+133.2}
\newcommand{\AlphaTwoDQN}{+83.1}
\newcommand{\AlphaTwoDQNCILow}{+64.5}
\newcommand{\AlphaTwoDQNCIHigh}{+99.1}
\newcommand{\AlphaTwoATwoC}{+153.6}
\newcommand{\AlphaTwoATwoCCILow}{+145.4}
\newcommand{\AlphaTwoATwoCCIHigh}{+161.4}
\newcommand{\AlphaTwoRF}{+113.2}
\newcommand{\AlphaTwoOracle}{+194.8}
\newcommand{\AlphaTwoGap}{+8.4}
\newcommand{\AlphaTwoVerdict}{overlapping}
\newcommand{\AlphaFourPPO}{+121.3}
\newcommand{\AlphaFourPPOCILow}{+108.6}
\newcommand{\AlphaFourPPOCIHigh}{+133.1}
\newcommand{\AlphaFourDQN}{+72.5}
\newcommand{\AlphaFourDQNCILow}{+53.5}
\newcommand{\AlphaFourDQNCIHigh}{+90.0}
\newcommand{\AlphaFourATwoC}{+138.7}
\newcommand{\AlphaFourATwoCCILow}{+128.0}
\newcommand{\AlphaFourATwoCCIHigh}{+150.0}
\newcommand{\AlphaFourRF}{+94.4}
\newcommand{\AlphaFourOracle}{+194.8}
\newcommand{\AlphaFourGap}{+26.9}
\newcommand{\AlphaFourVerdict}{significant}
\newcommand{\AlphaSixPPO}{+113.3}
\newcommand{\AlphaSixPPOCILow}{+100.4}
\newcommand{\AlphaSixPPOCIHigh}{+126.1}
\newcommand{\AlphaSixDQN}{+77.4}
\newcommand{\AlphaSixDQNCILow}{+60.4}
\newcommand{\AlphaSixDQNCIHigh}{+93.0}
\newcommand{\AlphaSixATwoC}{+151.4}
\newcommand{\AlphaSixATwoCCILow}{+141.9}
\newcommand{\AlphaSixATwoCCIHigh}{+159.8}
\newcommand{\AlphaSixRF}{+64.0}
\newcommand{\AlphaSixOracle}{+194.8}
\newcommand{\AlphaSixGap}{+49.3}
\newcommand{\AlphaSixVerdict}{significant}
\newcommand{\AlphaEightPPO}{+135.2}
\newcommand{\AlphaEightPPOCILow}{+124.9}
\newcommand{\AlphaEightPPOCIHigh}{+145.1}
\newcommand{\AlphaEightDQN}{+75.7}
\newcommand{\AlphaEightDQNCILow}{+56.3}
\newcommand{\AlphaEightDQNCIHigh}{+93.0}
\newcommand{\AlphaEightATwoC}{+137.5}
\newcommand{\AlphaEightATwoCCILow}{+126.5}
\newcommand{\AlphaEightATwoCCIHigh}{+148.0}
\newcommand{\AlphaEightRF}{+20.5}
\newcommand{\AlphaEightOracle}{+194.8}
\newcommand{\AlphaEightGap}{+114.7}
\newcommand{\AlphaEightVerdict}{significant}
\newcommand{\AlphaOnePPO}{+131.9}
\newcommand{\AlphaOnePPOCILow}{+120.7}
\newcommand{\AlphaOnePPOCIHigh}{+141.9}
\newcommand{\AlphaOneDQN}{+67.6}
\newcommand{\AlphaOneDQNCILow}{+49.6}
\newcommand{\AlphaOneDQNCIHigh}{+85.0}
\newcommand{\AlphaOneATwoC}{+142.0}
\newcommand{\AlphaOneATwoCCILow}{+130.7}
\newcommand{\AlphaOneATwoCCIHigh}{+153.4}
\newcommand{\AlphaOneRF}{-29.3}
\newcommand{\AlphaOneOracle}{+194.8}
\newcommand{\AlphaOneGap}{+161.2}
\newcommand{\AlphaOneVerdict}{significant}
\newcommand{\HeadlineAlpha}{0.4}
\newcommand{\AnchorPPO}{+138.6}
\newcommand{\AnchorRF}{+136.5}
\newcommand{\OracleCeiling}{+194.8}
\newcommand{\CouplingGapCoupled}{-63.1}
\newcommand{\CouplingGapOutcome}{-63.0}
\newcommand{\CouplingBestCoupled}{DQN}
\newcommand{\CouplingBestOutcome}{A2C}
\newcommand{\CouplingDQNCoupled}{+226.2}
\newcommand{\CouplingPPOCoupled}{+162.4}
\newcommand{\CouplingATwoCCoupled}{+144.8}
\newcommand{\CouplingDQNOutcome}{-8.6}
\newcommand{\CouplingPPOOutcome}{+126.2}
\newcommand{\CouplingATwoCOutcome}{+146.1}
% --- F10 environment-difficulty sweep (PPO/A2C/DQN) ---
\newcommand{\FtenZeroPPO}{-46.2}
\newcommand{\FtenZeroPPOCILow}{-62.4}
\newcommand{\FtenZeroPPOCIHigh}{-32.2}
\newcommand{\FtenTwoPPO}{+5.1}
\newcommand{\FtenTwoPPOCILow}{-12.4}
\newcommand{\FtenTwoPPOCIHigh}{+23.0}
\newcommand{\FtenFourPPO}{+69.6}
\newcommand{\FtenFourPPOCILow}{+48.1}
\newcommand{\FtenFourPPOCIHigh}{+93.0}
\newcommand{\FtenSixPPO}{+90.9}
\newcommand{\FtenSixPPOCILow}{+74.0}
\newcommand{\FtenSixPPOCIHigh}{+107.0}
\newcommand{\FtenEightPPO}{+104.1}
\newcommand{\FtenEightPPOCILow}{+84.4}
\newcommand{\FtenEightPPOCIHigh}{+123.0}
\newcommand{\FtenOnePPO}{+142.6}
\newcommand{\FtenOnePPOCILow}{+132.6}
\newcommand{\FtenOnePPOCIHigh}{+152.2}
\newcommand{\FtenZeroDQN}{-79.6}
\newcommand{\FtenZeroDQNCILow}{-118.2}
\newcommand{\FtenZeroDQNCIHigh}{-43.1}
\newcommand{\FtenTwoDQN}{-36.8}
\newcommand{\FtenTwoDQNCILow}{-73.0}
\newcommand{\FtenTwoDQNCIHigh}{+0.7}
\newcommand{\FtenFourDQN}{-10.3}
\newcommand{\FtenFourDQNCILow}{-39.6}
\newcommand{\FtenFourDQNCIHigh}{+20.5}
\newcommand{\FtenSixDQN}{+44.6}
\newcommand{\FtenSixDQNCILow}{+10.4}
\newcommand{\FtenSixDQNCIHigh}{+82.3}
\newcommand{\FtenEightDQN}{+65.4}
\newcommand{\FtenEightDQNCILow}{+37.4}
\newcommand{\FtenEightDQNCIHigh}{+94.8}
\newcommand{\FtenOneDQN}{+78.3}
\newcommand{\FtenOneDQNCILow}{+41.8}
\newcommand{\FtenOneDQNCIHigh}{+110.3}
\newcommand{\FtenZeroATwoC}{-7.2}
\newcommand{\FtenZeroATwoCCILow}{-21.0}
\newcommand{\FtenZeroATwoCCIHigh}{+5.4}
\newcommand{\FtenTwoATwoC}{+25.9}
\newcommand{\FtenTwoATwoCCILow}{+18.1}
\newcommand{\FtenTwoATwoCCIHigh}{+33.2}
\newcommand{\FtenFourATwoC}{+72.8}
\newcommand{\FtenFourATwoCCILow}{+53.6}
\newcommand{\FtenFourATwoCCIHigh}{+88.0}
\newcommand{\FtenSixATwoC}{+101.5}
\newcommand{\FtenSixATwoCCILow}{+83.6}
\newcommand{\FtenSixATwoCCIHigh}{+116.7}
\newcommand{\FtenEightATwoC}{+123.0}
\newcommand{\FtenEightATwoCCILow}{+108.0}
\newcommand{\FtenEightATwoCCIHigh}{+139.0}
\newcommand{\FtenOneATwoC}{+147.3}
\newcommand{\FtenOneATwoCCILow}{+136.5}
\newcommand{\FtenOneATwoCCIHigh}{+156.5}
\newcommand{\FtenGSevenThreePPO}{Yes}
\newcommand{\FtenGSevenThreeDQN}{Yes}
\newcommand{\FtenGSevenThreeATwoC}{Yes}
% --- F17 evasion-before-commit sweep (PPO/A2C/DQN) ---
\newcommand{\FseventeenZeroReward}{+123.2}
\newcommand{\FseventeenZeroCompromise}{0.451}
\newcommand{\FseventeenTwentyFiveReward}{+114.3}
\newcommand{\FseventeenTwentyFiveCompromise}{0.490}
\newcommand{\FseventeenFiftyReward}{+105.4}
\newcommand{\FseventeenFiftyCompromise}{0.534}
\newcommand{\FseventeenSeventyFiveReward}{+91.9}
\newcommand{\FseventeenSeventyFiveCompromise}{0.556}
\newcommand{\FseventeenZeroPPOReward}{+123.2}
\newcommand{\FseventeenZeroPPOCompromise}{0.451}
\newcommand{\FseventeenTwentyFivePPOReward}{+114.3}
\newcommand{\FseventeenTwentyFivePPOCompromise}{0.490}
\newcommand{\FseventeenFiftyPPOReward}{+105.4}
\newcommand{\FseventeenFiftyPPOCompromise}{0.534}
\newcommand{\FseventeenSeventyFivePPOReward}{+91.9}
\newcommand{\FseventeenSeventyFivePPOCompromise}{0.556}
\newcommand{\FseventeenZeroDQNReward}{+68.7}
\newcommand{\FseventeenZeroDQNCompromise}{0.677}
\newcommand{\FseventeenTwentyFiveDQNReward}{+57.8}
\newcommand{\FseventeenTwentyFiveDQNCompromise}{0.689}
\newcommand{\FseventeenFiftyDQNReward}{+56.4}
\newcommand{\FseventeenFiftyDQNCompromise}{0.702}
\newcommand{\FseventeenSeventyFiveDQNReward}{+50.5}
\newcommand{\FseventeenSeventyFiveDQNCompromise}{0.706}
\newcommand{\FseventeenZeroATwoCReward}{+142.6}
\newcommand{\FseventeenZeroATwoCCompromise}{0.233}
\newcommand{\FseventeenTwentyFiveATwoCReward}{+135.5}
\newcommand{\FseventeenTwentyFiveATwoCCompromise}{0.311}
\newcommand{\FseventeenFiftyATwoCReward}{+124.8}
\newcommand{\FseventeenFiftyATwoCCompromise}{0.386}
\newcommand{\FseventeenSeventyFiveATwoCReward}{+112.7}
\newcommand{\FseventeenSeventyFiveATwoCCompromise}{0.412}
\newcommand{\FseventeenCILowDegradation}{33.9}
\newcommand{\FseventeenPPOCILowDegradation}{33.9}
\newcommand{\FseventeenGSevenTenPPO}{No}
\newcommand{\FseventeenDQNCILowDegradation}{16.8}
\newcommand{\FseventeenGSevenTenDQN}{Yes}
\newcommand{\FseventeenATwoCCILowDegradation}{28.8}
\newcommand{\FseventeenGSevenTenATwoC}{Yes}
% --- F15 OOD-robustness (detector-independence) ---
\newcommand{\FfifteenBestRLPreventionLow}{0.71}
\newcommand{\FfifteenBestRLPreventionHigh}{0.85}
\newcommand{\FfifteenPPOPreventionLow}{0.71}
\newcommand{\FfifteenPPOPreventionHigh}{0.85}
\newcommand{\FfifteenRFPreventionLow}{0.00}
\newcommand{\FfifteenRFPreventionHigh}{0.15}
\newcommand{\FfifteenAdvantageLow}{+0.70}
\newcommand{\FfifteenAdvantageHigh}{+0.78}
\newcommand{\FfifteenSpearmanRho}{0.22}
\newcommand{\FfifteenSpearmanP}{0.54}
\newcommand{\FfifteenPearsonR}{-0.02}
\newcommand{\FfifteenPearsonP}{0.95}
\newcommand{\FfifteenOLSSlopeCILow}{-0.08}
\newcommand{\FfifteenOLSSlopeCIHigh}{+0.04}
\newcommand{\FfifteenSynFloodRecall}{0.998}
\newcommand{\FfifteenSynFloodAdvantage}{+0.73}
\newcommand{\FfifteenVulnScanRecall}{0.224}
\newcommand{\FfifteenVulnScanAdvantage}{+0.70}
\newcommand{\FfifteenDNSSpoofingRecall}{0.201}
% --- tuned RF-Acting stage-detector macro-F1 ---
\newcommand{\DetectorRfFoneBalanced}{0.924}
\newcommand{\NumTests}{462}

\input{generated/tables}
% ---

\renewcommand{\backrefpagesname}{Citado na(s) p\'{a}gina(s):~}
% Texto padr\~{a}o antes do n\'{u}mero das p\'{a}ginas
\renewcommand{\backref}{}
% Define os textos da cita\c{c}\~{a}o
\renewcommand*{\backrefalt}[4]{
	\ifcase #1 %
		Nenhuma cita\c{c}\~{a}o no texto.%
	\or
		Citado na p\'{a}gina #2.%
	\else
		Citado #1 vezes nas p\'{a}ginas #2.%
	\fi}%
% ---


% ---
% Configura\c{c}\~{o}es de apar\^{e}ncia do PDF final

% alterando o aspecto da cor azul
\definecolor{blue}{RGB}{41,5,195}

% informa\c{c}\~{o}es do PDF
\makeatletter
\hypersetup{
        %pagebackref=true,
		pdftitle={A Kill-Chain-Aware Reinforcement Learning Defense Framework for IoT Networks},
		pdfauthor={Felipe Augusto Oliveira dos Santos},
    pdfsubject={Master's dissertation on detector-free reinforcement learning for kill-chain-aware IoT defense},
	  pdfcreator={LaTeX with abnTeX2},
		pdfkeywords={IoT security; deep reinforcement learning; cyber kill chain; partial observability; CICIoT2023; reward specification; reproducibility},
		hidelinks,					% desabilita as bordas dos links
		colorlinks=false,       	% false: boxed links; true: colored links
    linkcolor=blue,          	% color of internal links
    citecolor=blue,        		% color of links to bibliography
    filecolor=magenta,      	% color of file links
		urlcolor=blue,
%		linkbordercolor={1 1 1},	% set to white
		bookmarksdepth=4
}
\makeatother
% ---

% ---
% Altera as margens padrões
% ---
\setlrmarginsandblock{3cm}{2cm}{*}
\setulmarginsandblock{3cm}{2cm}{*}
\checkandfixthelayout
% ---

% ---
% Espa\c{c}amentos entre linhas e par\'{a}grafos
% ---
\linespread{1.3}

% O tamanho do par\'{a}grafo \'{e} dado por:
\setlength{\parindent}{2.0cm}

% Controle do espa\c{c}amento entre um par\'{a}grafo e outro:
\setlength{\parskip}{0.2cm}  % tente tamb\'{e}m \onelineskip

% ---
% Informacoes de dados para CAPA e FOLHA DE ROSTO
% ---
\titulo{A Kill-Chain-Aware Reinforcement Learning Defense Framework for IoT Networks\\[12pt]{\large Um Arcabou\c{c}o de Defesa Baseado em Aprendizado por Refor\c{c}o Ciente da Cadeia de Elimina\c{c}\~{a}o para Redes IoT}}
\autor{Felipe Augusto Oliveira dos Santos}
\local{Campinas}
\data{2026}
\orientador{Prof. Dr. Denis Fantinato}
\instituicao{%
    UNIVERSIDADE ESTADUAL DE CAMPINAS
    \par
    Faculdade de Engenharia El\'{e}trica e de Computa\c{c}\~{a}o
    }
%% Master's dissertation — final defense version (Tese de Mestrado - Defesa Final)
\tipotrabalho{Disserta\c{c}\~{a}o (Mestrado)}
\preambulo{Disserta\c{c}\~{a}o apresentada \`{a} Faculdade de Engenharia El\'{e}trica e de Computa\c{c}\~{a}o da Universidade Estadual de Campinas como parte dos requisitos exigidos para a obten\c{c}\~{a}o do t\'{\i}tulo de Mestre em Engenharia El\'{e}trica, na \'{A}rea de Engenharia de Computa\c{c}\~{a}o.}
% ---
